\section{The measure}\label{sec:measure}\label{sec:construct}

\subsection{Two readers, one living backwards}

Imagine two readers of a filing's goods/services description, each having
seen a disjoint half of its industry. One has read only what the industry
filed in the five years before it. The other, living backwards as Merlin
was said to, has read only the five years after. Each is asked how
surprising the wording is. The measurement is how much they disagree: a
description that startles the first reader and not the second arrived
before its industry's language moved that way.

Both windows are anchored on the filing's own date, spanning the 1,826
days before it and the 1,826 days after; both exclude the filing date
itself, and inside a window every day counts the same. The measure asks
whether wording is unusual against a period, not whether it is ahead of
a trend inside that period; it has no momentum channel.

\subsection{What is measured}

The quantity is information content in Shannon's sense. Something that
happens with probability $p$ carries $\log(1/p)$ of information when it
happens: a common thing carries little, a rare thing a lot. Here the
things are the themes a description draws on, and the probabilities are
how often the filing's own industry drew on each of them in the
reference window. A filing's score asks, in effect: how surprising is
this mix of themes to someone who has only seen what the industry files?

The apparatus is that of \citet{Murdock2017}, developed on Darwin's
reading notebooks, and \citet{Barron2018}, who extended it to French
Revolution parliamentary debates; they call the two quantities
\emph{novelty} and \emph{transience} and their signed difference
\emph{resonance}. The arithmetic carries over unchanged; the
interpretation does not, and Section~\ref{sec:related} renames the
quantities for a commercial record in which the author is usually
counsel. One usage note carries through everywhere: in those settings
latent Dirichlet allocation identifies \emph{topics} --- what a document
is about. A goods/services description is not about anything in that
sense; it lists components, and across millions of filers a cluster of
co-occurring terms (\emph{video, music, games, electronic, audio}) is a
component many offerings share. \emph{Theme} is the word used
throughout; \emph{topic} appears only in the model's name.

\subsection{From a description to two numbers}\label{sec:definitions}

\paragraph{Counting words.} A description is reduced to a bag of words:
lowercased, split into tokens of three or more letters, and counted,
with adjacent pairs counted as tokens in their own right so that
\emph{computer software} registers alongside \emph{computer} and
\emph{software}. Take the filing for \textsc{Amazon.com}, lodged in the
advertising and retail class on 18 April 1997:

\begin{quote}\small
computerized on line search and ordering service featuring the wholesale
and retail distribution of books, music, motion pictures, multimedia
products and computer software in the form of printed books,
audiocassettes, videocassettes, compact disks, floppy disks, CD ROMs,
and direct digital transmission
\end{quote}

\noindent Of its 45 distinct terms, 36 are in the theme model's
vocabulary, among them \emph{computerized}, \emph{search},
\emph{ordering service}, \emph{retail} and \emph{digital transmission}.

\paragraph{Grouping words into themes.} Each filing is then re-described
in a fixed, much smaller set of coordinates. We run latent Dirichlet
allocation once over a stratified sample of the whole corpus, all 45
classes together. The model develops a list of 50 themes --- clusters of
words that tend to occur together --- and infers, for any filing, to
what degree its description invokes each one. Nothing supervises this,
and a theme is read off from its leading words: one collects the
vocabulary of media and entertainment goods, another retail store
services, another downloadable and online software. This build --- fifty
global themes, per-filing five-year windows, flat weighting --- is the
production scoring, used for every estimate unless the text states
another. The themes are shared across classes; what is specific to a
class is the reference, the average theme proportions of that class's
own filings in the window. The Amazon filing draws on six themes at more
than a twentieth of its weight: $0.41$ on media and entertainment goods,
$0.15$ on computer and information services, $0.13$ on retail store
services, $0.10$ on printed matter, $0.09$ on downloadable and online
software, and $0.07$ on marketing of consumer goods. The online appendix
lists every theme with its leading words and its weight in each
class.\footnote{\onlineappendixnote} Fifty is coarse, roughly one theme
per industry; Section~\ref{sec:three_scorings} refits at 500 themes and
at 50 themes within each class, and the supplementary material carries a
resolution sweep and a second fit at fixed resolution with a different
seed.

\paragraph{Comparing against the two references.} Write $P$ for the
filing's distribution over themes, and $Q^{-}$ and $Q^{+}$ for that
distribution averaged over every same-class filing in the two windows.
How far $P$ sits from a reference $Q$ is the Kullback--Leibler
divergence,
\[
  \mathrm{KL}(P \,\|\, Q) \;=\; \sum_{k} P_k \log \frac{P_k}{Q_k},
\]
zero when the two distributions agree and growing as the filing puts
weight where the industry puts little, measured in nats. Write
$K^{-} = \mathrm{KL}(P \,\|\, Q^{-})$ and
$K^{+} = \mathrm{KL}(P \,\|\, Q^{+})$: the filing's past-facing and
future-facing surprise, what the first reader feels and the second.

\paragraph{Taking the difference.} Amazon's filing is scored against
41{,}166 earlier filings and 146{,}532 later ones, giving
$K^{-} = 1.286$ and $K^{+} = 1.163$. The two axes used throughout are
their average and their signed difference:
\[
  \underbrace{A \;=\; \tfrac{1}{2}\left(K^{-} + K^{+}\right)}_{\text{atypicality}},
  \qquad
  \underbrace{L \;=\; K^{-} - K^{+}}_{\text{lead}} .
\]
For Amazon, $A = 1.225$ and $L = +0.123$: against its own class, its
atypicality is at the 47th percentile and its lead at the 94th. The
industry moved toward this description in the five years after it was
filed and had not been writing that way in the five years before; the
three themes that carry most of its description all became more common
in the class over the following five years.

Why the average and the difference, rather than the two surprises
themselves? Because $K^{-}$ and $K^{+}$ correlate at $0.988$: a filing
far from its class's past is almost always far from its future too, so
entering both is close to entering one variable twice. The average
collects that large shared component and the difference collects the
small residual that carries the timing; nothing more than a sum and a
difference is taken, with no weights estimated, and the two resulting
regressors are nearly uncorrelated ($+0.036$). The cost is that $L$ is a
small residual of two large quantities, about a sixth of either level's
spread, which the supplementary material quantifies.

\paragraph{Why the words are not scored directly.} The words themselves
could be compared to the reference, and the paper reports that scoring
as a cross-check. It cannot be the primary one, because a divergence
measured on words mostly counts the length of the filing. For any two
distributions,
\[
  \mathrm{KL}(P \,\|\, Q) \;=\; \underbrace{-\textstyle\sum_k P_k \log Q_k}_{\text{how unexpected these words are}}
  \;-\; \underbrace{\left(-\textstyle\sum_k P_k \log P_k\right)}_{\text{how spread out the filing is}} ,
\]
and the spread of a filing over words is bounded by its distinct-term
count: a description using $k$ terms cannot spread past $\log k$, while
the reference spreads over about 900 terms. A five-word description is
charged roughly five nats of width before any content is considered. In
3{,}000 long software-class filings cut down to random handfuls of words
and rescored, a fiftyfold cut in length moves the surprise term by
$0.016$ nats and the spread term by $2.9$: the score change is the
length, not the content. Term-scored atypicality correlates with log
distinct terms at $-0.68$ in the software class. Themes break the
coupling because the 50 coordinates are always occupied; theme-scored
atypicality correlates with log term count at $-0.22$ there, and
$-0.10$ in advertising and retail. What term scoring alone resolves is
novel word \emph{combinations} --- of Amazon's inventive bigrams,
nothing survives into any theme basis --- so theme scoring carries every
estimate and term scoring every worked example (Supplementary Material,
\S S2, traces three filings through both representations).

\paragraph{What the themes see, and what they do not.} The score reads
which themes a filing draws on and how heavily, not how it combines
them: on a 150{,}000-filing software-class sample, the rarity of a
filing's two leading themes explains 70\% of the variance in
atypicality and the unusualness of the pairing adds under two points. A
filing assembled from ordinary parts in a genuinely new configuration is
close to invisible to this measure. Scoring arrangements directly ---
a pair-level lead, and word-level recombination in the sense of
\citet{Fleming2001, Uzzi2013} --- shows the two channels acting in
opposite directions in the software class and neither generalizing
corpus-wide, so the construct's blindness to combination is a bounded
omission rather than a hidden bias in one direction (construction and
estimates in the supplementary material).

\subsection{Within class, not across it}

Goods/services descriptions are drafted to secure scope of protection,
and counsel reuse language from the USPTO's own Acceptable
Identification Manual, from house precedent, and from competitors'
successful filings. Identical wording across unrelated firms is
therefore common: 17\% of scored registrations share a value with
another filing in their class and registration year. Within a class, the
measure picks up departure from a shared drafting convention. Across
classes it does not: a class where the Manual supplies dense standard
language shows a compressed distribution of atypicality, and one where
firms describe novel services in prose a wide one, for reasons unrelated
to how innovative either industry is. Every estimate is therefore taken
within class and year, and magnitudes compare across classes only as
signs and orderings. The vocabulary is positional rather than causal: a
leading filing was ahead of its industry's language, which implies a
trajectory without asserting the filing caused it.
